\documentclass[14pt,a4paper]{extarticle}

% ---------- Кодировки, язык, шрифты ----------
% Поддержка русского текста при сборке pdfLaTeX без cmap/mmap.
% Рекомендуемый компилятор: pdfLaTeX, минимум два запуска для оглавления.
\usepackage[T2A]{fontenc}
\usepackage[utf8]{inputenc}
\usepackage[russian]{babel}

% ---------- Страница и интервалы (ГОСТ 7.32: A4, 14 пт, 1,5 интервала) ----------
\usepackage[a4paper,left=30mm,right=15mm,top=20mm,bottom=20mm]{geometry}
\usepackage{setspace}
\onehalfspacing
\usepackage{indentfirst}
\setlength{\parindent}{1.25cm}
\setlength{\parskip}{0pt}
\emergencystretch=2em

% ---------- Математика, таблицы, списки ----------
\usepackage{amsmath,amssymb}
\usepackage{array}
\usepackage{enumitem}
\setlist[itemize]{label=--, leftmargin=1.9cm, itemsep=0pt, topsep=2pt, parsep=0pt}

% ---------- Заголовки ----------
\usepackage{titlesec}
\titleformat{\section}{\normalfont\large\bfseries}{\thesection}{0.7em}{}
\titleformat{\subsection}{\normalfont\normalsize\bfseries}{\thesubsection}{0.7em}{}
\titlespacing*{\section}{0pt}{0pt}{14pt}
\titlespacing*{\subsection}{0pt}{12pt}{6pt}
% каждый раздел начинается с новой страницы
\newcommand{\sectionbreak}{\clearpage}
% разделы без номера, попадающие в содержание
\newcommand{\unnumsection}[1]{\section*{#1}\phantomsection\addcontentsline{toc}{section}{#1}}
\renewcommand{\contentsname}{Содержание}
\setcounter{tocdepth}{2}

% ---------- Подписи таблиц и листингов ----------
\usepackage{caption}
\DeclareCaptionLabelSeparator{emdash}{ --- }
\captionsetup{labelsep=emdash,justification=centering,font=normalsize,singlelinecheck=false}
\DeclareCaptionType{listing}[Листинг][Список листингов]
\captionsetup[table]{skip=6pt}
\captionsetup[listing]{skip=4pt,hypcap=false}

% ---------- Листинги ----------
% Стандартное окружение verbatim встроено в LaTeX, пакет не нужен.

% ---------- Ссылки ----------
\usepackage{xcolor}
\usepackage[unicode,colorlinks=true,linkcolor=black,citecolor=black,urlcolor=blue!40!black,
  pdftitle={Курсовая работа: бот-игрок в «Крестики-нолики»},
  pdfauthor={Переверзева М.А.}]{hyperref}
\DeclareUrlCommand\code{\urlstyle{tt}}
\renewcommand{\arraystretch}{1.25}

% ---------- Данные титульного листа (править здесь) ----------
\newcommand{\StudentName}{Переверзева М.А.}
\newcommand{\StudentGroup}{5130201/50301}
\newcommand{\SupervisorPost}{мл. науч. сотр., ассистент}
\newcommand{\SupervisorName}{Кирпиченко С.Р.}
% пометка для полей, которые нужно заполнить вручную
\newcommand{\todo}[1]{\textbf{[#1]}}

\begin{document}

% ======================= ТИТУЛЬНЫЙ ЛИСТ =======================
\begin{titlepage}
\thispagestyle{empty}
\singlespacing
\centering
{\large МИНИСТЕРСТВО НАУКИ И ВЫСШЕГО ОБРАЗОВАНИЯ\\РОССИЙСКОЙ ФЕДЕРАЦИИ\par}
\vspace{1.5ex}
{\normalsize федеральное государственное автономное образовательное учреждение высшего образования «Санкт-Петербургский политехнический университет Петра Великого»\par}
\vspace{2ex}
{\normalsize Институт компьютерных наук и кибербезопасности\\Высшая школа технологий искусственного интеллекта\par}
\vspace{2ex}
{\normalsize Направление: 02.03.01 Математика и компьютерные науки\par}
\vspace{4cm}
{\large КУРСОВАЯ РАБОТА\par}
\vspace{1.5ex}
{\normalsize по дисциплине «Программирование и алгоритмизация»\par}
\vspace{1.5ex}
{\normalsize Тема: бот-игрок в игре «Крестики-нолики» на основе минимакса\\с альфа-бета отсечением\par}
\vfill
\hfill\begin{minipage}{0.46\textwidth}
\raggedleft
Выполнила: студентка группы\\\StudentGroup\par
\vspace{1.2ex}
\rule{0.85\linewidth}{0.4pt}\\
\StudentName\par
\vspace{1.5ex}
Проверил: \SupervisorPost\par
\vspace{1.2ex}
\rule{0.85\linewidth}{0.4pt}\\
\SupervisorName
\end{minipage}
\vfill
Санкт-Петербург\\2026 г.
\end{titlepage}
\setcounter{page}{2}% титульный лист считается первой страницей

% ======================= СОДЕРЖАНИЕ =======================
\tableofcontents

% ======================= ВВЕДЕНИЕ =======================
\unnumsection{Введение}

Игра «крестики-нолики» --- классический пример дискретной игры с полной информацией, на примере которой удобно изучать методы поиска в пространстве состояний: полный перебор, минимакс, альфа-бета отсечение, эвристическую оценку позиции~\cite{russell}. Несмотря на простоту правил, вариант игры на большом поле ($20\times 20$) с препятствиями и увеличенной длиной выигрышной комбинации (5 в ряд) практически неотличим по сложности перебора от игры «гомоку» (five-in-a-row), для которой полный перебор невозможен даже на современных вычислительных мощностях: число возможных партий астрономически велико, а требование уложиться в 50--100~мс на ход делает задачу вычислительно нетривиальной.

Работа с подобными задачами полезна как учебное упражнение по алгоритмизации: она требует одновременно применения математической модели (оценочной функции, отсекающего поиска) и практических приёмов программирования (эффективных структур данных, инкрементального пересчёта, ограничения по времени). Аналогичные подходы используются в реальных игровых движках (шахматные и шашечные программы) и, в упрощённом виде, в системах принятия решений с ограниченным временем реакции.

В данной работе рассматривается разработка бота-игрока для крестиков-ноликов до пяти в ряд на поле $20\times 20$ с препятствиями, встраиваемого в готовую инфраструктуру игры (ядро правил, генератор препятствий, интерфейс игрока), предоставленную в рамках курса «Программирование и алгоритмизация»~\cite{repo}. Разработка велась с применением трекера задач и Kanban-доски (GitHub Projects~\cite{ghprojects}): курсовая работа разбита на отдельные задачи (Issues), каждая реализована в собственной ветке репозитория и присоединена к основной линии разработки через Pull Request, что отражено в истории версий проекта.

% ======================= 1. ПОСТАНОВКА ЗАДАЧИ =======================
\section{Постановка задачи}

Дано готовое ядро игры «крестики-нолики» (библиотека \code{tttcore})~\cite{repo}, реализующее:
\begin{itemize}
  \item хранение состояния игрового поля и правил перехода между состояниями (класс \code{State});
  \item генерацию препятствий на поле по параметрам \code{playable_part}, \code{max_obstacle_len}, \code{gap};
  \item игровой цикл, поочерёдно опрашивающий двух игроков через интерфейс \code{IPlayer}, до завершения партии;
  \item предкомпилированную закрытую реализацию двух игроков базового уровня («лёгкого» и «сложного») для сравнения.
\end{itemize}

Требуется реализовать собственного игрока (класс, реализующий интерфейс \code{IPlayer}), удовлетворяющего следующим условиям:
\begin{itemize}
  \item корректно играет по правилам --- никогда не выбирает занятую клетку или клетку вне поля;
  \item рассчитывает каждый ход не дольше 50--100~мс на рабочей станции уровня ноутбука;
  \item выигрывает у игрока базового уровня сложности более чем в 50\% партий (без учёта ничьих) и у игрока повышенного уровня сложности --- более чем в 25\% партий;
  \item работает на фиксированных параметрах партии: поле $20\times 20$, длина выигрышной линии 5, \code{playable_part} $=0{,}75$, \code{gap} $=1$, \code{max_obstacle_len} $=50$.
\end{itemize}

Дополнительно необходимо: словесно и математически описать применяемый алгоритм без привязки к терминам языка программирования (раздел~\ref{sec:math}), реализовать его на C++ в виде класса, встраиваемого в проект (раздел~\ref{sec:impl}), протестировать полученного бота и измерить время расчёта хода (раздел~\ref{sec:results}), подготовить отчёт о проделанной работе.

Результатом работы являются: исходный код бота (см. приложение~А), набор автоматических тестов на основе CTest~\cite{ctest}, конфигурация непрерывной интеграции (CI) на базе GitHub Actions~\cite{ghactions} и настоящий отчёт.

% ======================= 2. МАТЕМАТИЧЕСКОЕ ОПИСАНИЕ =======================
\section{Математическое описание}\label{sec:math}

\subsection{Модель игры}

Поле --- прямоугольная таблица размером 20 строк на 20 столбцов. Каждая клетка поля находится в одном из четырёх состояний: свободна, содержит препятствие, содержит знак первого игрока (X) или знак второго игрока (O). Игроки ходят по очереди в свободную клетку, начинает первый игрок. Партия завершается, когда какой-либо игрок составил линию из пяти или более своих знаков подряд по горизонтали, вертикали или одной из двух диагоналей, либо когда все клетки поля заняты.

Особое правило: если линию первым составил второй игрок, он побеждает немедленно. Если линию первым составил первый игрок, второй игрок получает право на последний ход; если он тоже успевает составить линию, объявляется ничья, иначе побеждает первый игрок.

\emph{Окном} будем называть любые пять клеток поля, идущих подряд вдоль строки, столбца или диагонали (под $45^\circ$). Составить выигрышную линию можно только внутри окна, свободного от препятствий и не содержащего знаков соперника.

\subsection{Оценка позиции}

Пусть $E$ --- число, характеризующее позицию с точки зрения бота: чем больше $E$, тем позиция выгоднее для бота, чем меньше (отрицательнее) --- тем выгоднее для соперника.

Величина $E$ равна сумме вкладов $v(W)$ по всем окнам $W$ поля:
\begin{equation}\label{eq:E}
E=\sum_{W} v(W).
\end{equation}
Вклад одного окна определяется так:
\begin{equation}\label{eq:v}
v(W)=\begin{cases}
0, & \text{\footnotesize если в окне есть препятствие, знаки обоих игроков или нет знаков;}\\
+w(k), & \text{\footnotesize если в окне $k$ знаков бота и нет знаков соперника;}\\
-w(k), & \text{\footnotesize если в окне $k$ знаков соперника и нет знаков бота.}
\end{cases}
\end{equation}
Здесь $k$ --- число знаков в окне, а весовая функция $w(k)$ задана в таблице~\ref{tab:w}.

\begin{table}[h!]
\captionsetup{position=above}
\caption{Весовая функция окна $w(k)$}\label{tab:w}
\centering\small
\begin{tabular}{|c|c|c|c|c|c|c|}
\hline
$k$ (число одинаковых знаков в окне) & 0 & 1 & 2 & 3 & 4 & 5\\ \hline
$w(k)$ & 0 & 1 & 10 & 300 & $20\,000$ & $10\,000\,000$\\ \hline
\end{tabular}
\end{table}

Значения $w(k)$ подобраны так, чтобы каждая дополнительная фигура в окне резко увеличивала вклад: наличие четырёх своих знаков в свободном окне (готовая угроза выигрыша следующим ходом) весит на четыре порядка больше, чем одиночный знак, а собранная линия из пяти оценивается как безусловная победа.

\subsection{Отбор ходов-кандидатов}\label{sec:m-cand}

Перебор всех свободных клеток поля на каждом уровне поиска приводит к неприемлемому времени работы, поэтому рассматриваются только клетки-кандидаты: свободные клетки, отстоящие не более чем на 2 клетки (по каждой координате) от какой-либо уже занятой клетки. Если поле полностью пустое (первый ход партии), кандидатом считается единственная клетка --- ближайшая к центру поля.

Каждому кандидату $x$ присваивается приоритет --- оценка немедленной полезности хода в эту клетку для игрока, который сейчас ходит (назовём его «текущим»). Приоритет вычисляется как сумма вкладов по всем окнам, содержащим клетку $x$ и свободным от препятствий:
\begin{itemize}
  \item пустое окно (в нём ещё нет ни одного знака) даёт вклад $+1$;
  \item окно, где после гипотетического хода будет $k$ знаков текущего игрока и ни одного знака соперника, даёт вклад $a(k)$ --- «атакующая» ценность;
  \item окно, где есть $k$ знаков соперника и ни одного знака текущего игрока, даёт вклад $d(k)$ --- «защитная» ценность, отражающая необходимость блокировать угрозу соперника.
\end{itemize}
Коэффициенты $a(k)$ и $d(k)$ приведены в таблице~\ref{tab:ad}.

\begin{table}[h!]
\captionsetup{position=above}
\caption{Весовые коэффициенты приоритета клетки}\label{tab:ad}
\centering\small
\begin{tabular}{|l|c|c|c|c|}
\hline
$k$ & 1 & 2 & 3 & 4\\ \hline
$a(k)$ --- свои знаки & 10 & 100 & $1\,000$ & $10\,000\,000$\\ \hline
$d(k)$ --- чужие знаки & 5 & 50 & 500 & $1\,000\,000$\\ \hline
\end{tabular}
\end{table}

Кандидаты сортируются по убыванию приоритета; в дальнейшем переборе участвуют только 12 лучших (на первом уровне перебора --- 30 лучших, поскольку там цена ошибки выше).

\emph{Форсированный ход.} Если у лучшего кандидата приоритет не меньше $1\,000\,000$ (это возможно только когда ход завершает линию из пяти своих знаков или блокирует линию соперника из четырёх), то остальные кандидаты не рассматриваются: форсированный ход заведомо не хуже любого другого.

\subsection{Поиск с альфа-бета отсечением}\label{sec:m-search}

Значение позиции для игрока, который сейчас должен ходить, определяется рекурсивно (алгоритм минимакс в форме негамакс~\cite{russell}: оценка всегда берётся с точки зрения ходящего, поэтому знак значения меняется на противоположный при переходе хода к сопернику):
\begin{itemize}
  \item если ход завершает линию из пяти и более своих знаков --- значение хода равно очень большому числу ($100\,000\,000$) минус номер полухода партии: чем быстрее достигается выигрыш, тем выше значение;
  \item если достигнута заданная глубина перебора $D$ полуходов и партия не завершена --- значением листа считается $E$, если сейчас ходит бот, и $-E$, если ходит соперник;
  \item иначе значение хода равно минус значению позиции для соперника, который делает следующий ход из получившегося состояния (рекурсивный вызов той же процедуры с противоположным знаком).
\end{itemize}
Значением позиции для игрока считается максимум значений всех его ходов-кандидатов.

\emph{Отсечение (альфа-бета).} Пусть при переборе ходов игрока уже найден ход со значением $\alpha$ (нижняя граница того, что игрок гарантированно может получить). При разборе очередного хода, как только у соперника находится ответ, после которого итоговое значение для игрока становится не выше $\alpha$, дальнейшие ответы соперника на этот ход можно не проверять: этот ход заведомо не лучше уже найденного, и разбор оставшихся ответов не может изменить общий результат перебора. Это позволяет отбрасывать значительную часть дерева перебора без потери точности итогового результата.

\subsection{Итеративное увеличение глубины и ограничение по времени}

Поскольку заранее не известно, какая глубина перебора уложится в отведённое время, применяется следующая схема:
\begin{itemize}
  \item перебор выполняется последовательно для глубин $D=2,4,6,\dots$ полуходов (чётные глубины выбраны так, чтобы оценка листа всегда ставилась после ответного хода соперника, а не посреди обмена ходами);
  \item перед каждой новой глубиной лучший ход, найденный на предыдущей глубине, рассматривается первым --- это ускоряет отсечения на следующей итерации;
  \item для каждого хода отслеживается прошедшее время; как только оно превышает установленный лимит (по умолчанию 70~мс), текущая, ещё не завершённая глубина прекращается, и используется лучший ход последней полностью завершённой глубины;
  \item если на какой-либо глубине найден гарантированный выигрыш или неизбежный проигрыш, дальнейшее увеличение глубины прекращается, так как результат уже не может измениться.
\end{itemize}
Итоговый ход --- координаты клетки, выбранной описанной процедурой на последней успешно завершённой (или прерванной по таймеру) глубине перебора.

% ======================= 3. ОСОБЕННОСТИ РЕАЛИЗАЦИИ =======================
\section{Особенности реализации}\label{sec:impl}

\subsection{Программный интерфейс}

Алгоритм реализован в виде класса \code{MyPlayer} (файлы \code{src/player/my_player.hpp}, \code{src/player/my_player.cpp}), реализующего интерфейс игры \code{ttt::game::IPlayer}. Публичный интерфейс класса приведён в листинге~\ref{lst:iface}.

\captionof{listing}{Публичный интерфейс класса \texttt{MyPlayer}}\label{lst:iface}
\begin{verbatim}
class MyPlayer : public IPlayer {
public:
  MyPlayer(const char *name) : m_sign(Sign::NONE), m_name(name) {}
  void set_sign(Sign sign) override;
  Point make_move(const State &game) override;
  const char *get_name() const override;

  // Настройка и статистика (для тестов и отчета).
  void set_time_limit_ms(int ms) { m_time_limit_ms = ms; }
  int last_depth() { return m_last_depth; }
  ...
};
\end{verbatim}

Игра вызывает \code{set_sign(Sign)} один раз в начале партии, сообщая, каким знаком играет бот, и \code{make_move(const State&)} перед каждым ходом бота, ожидая в ответ координаты свободной клетки. Дополнительные методы (\code{set_time_limit_ms}, \code{last_depth}, \code{last_nodes}, \code{last_time_ms}) предназначены для настройки лимита времени и получения статистики поиска в автоматических тестах и при измерении производительности.

\subsection{Внутреннее представление поля}

Состояние \code{State}, полученное от игры, копируется в собственный массив \code{std::vector<std::int8_t>} размером $\text{rows}\times\text{cols}$ (метод \code{load}), где каждой клетке соответствует одно из четырёх значений: свободна, знак бота, знак соперника, препятствие. Такое представление занимает один байт на клетку (в отличие от объектного \code{State}) и позволяет быстро делать и откатывать пробные ходы при переборе без обращения к интерфейсу игры.

Дополнительно хранится вспомогательный массив \code{m_near} того же размера, в котором для каждой клетки накапливается число занятых клеток в радиусе 2 вокруг неё (метод \code{bump_near}). Ненулевое значение в этом массиве --- признак того, что клетка является кандидатом на ход; это позволяет перечислять кандидатов за один проход по полю без вычисления расстояний в момент перебора.

\subsection{Инкрементальный пересчёт оценки}

Оценка позиции $E$ (формула~\eqref{eq:E}) при полной перезагрузке поля считается суммированием вкладов всех окон; листинг~\ref{lst:window} показывает вычисление вклада одного окна.

\captionof{listing}{Вычисление вклада одного окна в оценку позиции}\label{lst:window}
\begin{verbatim}
long long MyPlayer::window_value(int sx, int sy, int d) const {
  int me = 0, opp = 0;
  for (int i = 0; i < m_len; ++i) {
    std::int8_t c = at(sx + DX[d] * i, sy + DY[d] * i);
    if (c == WALL)
      return 0;
    if (c == ME)
      ++me;
    else if (c == OPP)
      ++opp;
  }
  if (me && opp)
    return 0;
  if (me)
    return m_W[me];
  if (opp)
    return -m_W[opp];
  return 0;
}
\end{verbatim}

Пересчитывать сумму по всем окнам поля после каждого пробного хода в переборе было бы слишком медленно. Поэтому при постановке или снятии знака в клетке (методы \code{apply}/\code{undo}) из общей суммы вычитаются старые вклады только тех окон, что проходят через изменённую клетку (таких окон не более $4\times 5=20$ --- по четыре направления и по пять положений окна), а затем добавляются новые вклады тех же окон. Такой инкрементальный пересчёт занимает время, не зависящее от размера поля, и является основным фактором, позволяющим уложиться в отведённый лимит времени на ход при переборе на глубину в несколько полуходов.

\subsection{Генерация кандидатов и форсированные ходы}

Функция \code{gen} (листинг~\ref{lst:gen}) перечисляет клетки-кандидаты (те, где \code{m_near} $>0$), присваивает им приоритет через вспомогательную функцию \code{cell_priority} (реализующую формулу из раздела~\ref{sec:m-cand}) и оставляет только \code{beam} лучших по приоритету, используя частичную сортировку \code{std::partial_sort} со сложностью $O(N\log \text{beam})$~\cite{cormen,cppref}. Если приоритет лучшего кандидата не ниже порога, соответствующего форсированному ходу, список кандидатов сокращается до одного элемента.

\captionof{listing}{Генерация и отбор кандидатов на ход}\label{lst:gen}
\begin{verbatim}
int MyPlayer::gen(int ply, std::int8_t mover, int beam) {
  auto &v = m_cands[ply];
  v.clear();
  for (int x = 0; x < m_nx; ++x)
    for (int y = 0; y < m_ny; ++y) {
      int idx = x * m_ny + y;
      if (m_b[idx] == EMPTY && m_near[idx] > 0)
        v.push_back({x, y, cell_priority(x, y, mover)});
    }
  int n = std::min<int>(beam, (int)v.size());
  std::partial_sort(v.begin(), v.begin() + n, v.end(),
                    [](const Cand &a, const Cand &b) { return a.pr > b.pr; });
  // Форсированные ходы: есть ход, собирающий линию, или надо блокировать
  // соперника, у которого до линии не хватает одной клетки.
  if (n > 0 && v[0].pr >= m_def[m_len - 1])
    n = 1;
  return n;
}
\end{verbatim}

\subsection{Поиск и управление временем}

Собственно перебор реализован методом \code{negamax} (рекурсивный обход дерева ходов с альфа-бета отсечением, раздел~\ref{sec:m-search}) и методом \code{search_root} (отдельная реализация верхнего уровня перебора, которая, в отличие от внутренних узлов, запоминает не только оценку, но и сами координаты лучшего хода). Итеративное углубление и контроль времени реализованы в \code{make_move}: перед каждым вызовом запоминается момент начала расчёта (\code{std::chrono::steady_clock}~\cite{cppref}), вычисляется дедлайн, и внутри \code{negamax} раз в 256 просмотренных узлов проверяется, не истекло ли отведённое время; при превышении лимита текущая ветвь перебора прерывается флагом \code{m_timeout}, и результат текущей (незавершённой) глубины отбрасывается.

Начальным запасным ходом, который используется, если не завершилась ни одна полная глубина перебора (например, на очень перегруженном препятствиями поле), служит кандидат с наибольшим приоритетом из \code{gen} --- то есть локально лучший ход по одноходовой эвристике.

\subsection{Используемые структуры данных и случайность}

Алгоритм детерминирован: генератор случайных чисел не используется --- при одинаковом входном состоянии поля бот всегда выбирает один и тот же ход. Из структур данных, помимо описанных выше плоских массивов на основе \code{std::vector}, используется небольшой массив \code{std::vector<Cand>} на каждый уровень перебора (переиспользуемый между вызовами, чтобы не выделять память многократно), где \code{Cand} --- простая структура «координаты + приоритет».

Все веса (\code{m_W}, \code{m_att}, \code{m_def}) --- константы, зависящие только от длины выигрышной линии (\code{win_len}) и вычисляемые один раз при загрузке состояния (метод \code{load}), что позволяет использовать один и тот же класс без изменений для других значений длины линии, если параметры игры изменятся.

\subsection{Организация разработки и непрерывная интеграция}\label{sec:ci}

Разработка велась по задачам (Issues) на Kanban-доске GitHub Projects~\cite{ghprojects} с тремя колонками (Backlog, In Progress, Done). Каждая задача реализована в отдельной ветке репозитория (\code{feature/#N_...}), коммиты помечены номером задачи (например, \code{[#3] Implement MyPlayer...}), а слияние с веткой \code{dev} выполнено через Pull Request с пометкой \code{Closes #N}, закрывающей соответствующий Issue.

Для автоматической проверки работоспособности проекта настроен процесс непрерывной интеграции на основе GitHub Actions~\cite{ghactions} (файл \code{.github/workflows/ci.yml}, листинг~\ref{lst:ci}). При каждом push в ветки \code{main}, \code{dev} и \code{feature/**} и при открытии Pull Request в \code{main} или \code{dev} запускается сборка проекта в конфигурации Release и полный прогон тестов через CTest.

\captionof{listing}{Конфигурация непрерывной интеграции (\texttt{.github/workflows/ci.yml})}\label{lst:ci}
\begin{verbatim}
# Кладётся в .github/workflows/ci.yml
# Собирает проект (без предкомпилированной библиотеки baseline, так как она
# закрытая и недоступна CI) и запускает CTest.
name: CI

on:
  push:
    branches: [ main, dev, "feature/**" ]
  pull_request:
    branches: [ main, dev ]

jobs:
  build-and-test:
    runs-on: ubuntu-latest
    steps:
      - name: Checkout
        uses: actions/checkout@v4

      - name: Configure
        run: |
          mkdir -p build
          cd build
          cmake .. -DCMAKE_BUILD_TYPE=Release

      - name: Build
        run: cmake --build build --config Release -j

      - name: Test
        run: |
          cd build
          ctest --output-on-failure
\end{verbatim}

% ======================= 4. РЕЗУЛЬТАТЫ =======================
\section{Результаты работы программы}\label{sec:results}

\subsection{Методика тестирования}

Все тесты проводились при фиксированных параметрах партии, установленных заданием: поле $20\times 20$, длина выигрышной линии 5, \code{playable_part} $=0{,}75$, \code{gap} $=1$, \code{max_obstacle_len} $=50$. Лимит времени на ход бота --- 70~мс (метод \code{set_time_limit_ms}), что оставляет запас до предельных 100~мс с учётом накладных расходов на копирование состояния.

Предоставленная библиотека игроков базового уровня (\code{libtttcore.a}) собрана под Linux (формат ELF) и не может быть подключена к проекту, собранному компилятором MSVC в Windows. Поэтому итоговые испытания выполнены в среде WSL (Ubuntu) на компьютере с ОС Windows: компилятор GCC 15.2.0, CMake 4.2.3~\cite{cmake}, конфигурация сборки Release. Процессор: \todo{укажите модель процессора}.

В проекте используются следующие автоматические тесты (папка \code{tests/}):
\begin{itemize}
  \item \code{smoke_test_player} (\code{tests/test_my_player.cpp}) --- базовая проверка работоспособности игрока, предоставленная шаблоном проекта;
  \item \code{test_player_stats} (\code{tests/test_stats.cpp}) --- 100 партий бота против самого себя, проверка устойчивости бота;
  \item \code{test_my_bot} (\code{tests/test_my_bot.cpp}, добавлен в рамках данной работы) --- три целевые проверки: корректность возвращаемых ходов (клетка внутри поля и свободна), полная партия бота против самого себя без дисквалификаций, время одного хода не более 100~мс;
  \item \code{test_with_baseline} (\code{tests/test_baseline.cpp}) --- одна партия против игрока базового уровня с выводом поля после каждого хода, используется для наглядной проверки;
  \item \code{test_stats_vs_baseline} (\code{tests/test_stats_vs_baseline.cpp}) --- серия из 100 партий против игрока базового уровня с подсчётом побед, поражений, ничьих и времени хода.
\end{itemize}

Для воспроизведения результатов в WSL выполняются команды \code{cmake .. -DBUILD_TTTCORE=PREBUILT -DCMAKE_BUILD_TYPE=Release} и \code{cmake --build . -j}, после чего запускается \code{./tests/test_stats_vs_baseline}. Для испытания против сильного игрока в этом файле заменяется строка получения соперника: вместо \code{ttt::baseline::get_easy_player} используется \code{ttt::baseline::get_harder_player}, остальной код не меняется. Результаты сохранялись в текстовые файлы командой \code{tee}.

Сборка и запуск тестов автоматизированы через GitHub Actions (раздел~\ref{sec:ci}): при каждом коммите в отслеживаемые ветки CI собирает проект в конфигурации Release и выполняет набор тестов командой \code{ctest --output-on-failure}. В CI предкомпилированная библиотека недоступна, поэтому тесты против baseline там не выполняются. Прогон CI на ветке \code{dev} завершился успешно (все доступные тесты пройдены, общее время около 24 минут).

\subsection{Результаты автоматических тестов}

\begin{table}[h!]
\captionsetup{position=above}
\caption{Результаты прогона тестов CTest}\label{tab:ctest}
\centering\small
\begin{tabular}{|>{\raggedright\arraybackslash}p{4.6cm}|>{\raggedright\arraybackslash}p{6.6cm}|>{\raggedright\arraybackslash}p{3.8cm}|}
\hline
Тест & Что проверяет & Результат\\ \hline
\code{smoke_test_player} & Базовая работоспособность игрока & Пройден\\ \hline
\code{test_my_bot} & Легальность хода, партия без дисквалификации, время хода $\le 100$~мс (3 проверки) & Пройден (3/3)\\ \hline
\code{test_player_stats} & 100 партий бота против самого себя & Пройден (в CI)\\ \hline
\code{test_with_baseline} & Одна партия против сильного игрока базового уровня & Завершена, ничья\\ \hline
\end{tabular}
\end{table}

\subsection{Результаты партий против базового уровня}

Серия из 100 партий против каждого из игроков базового уровня сыграна тестом \code{test_stats_vs_baseline}. Бот играл за первого игрока (X). Процент побед без учёта ничьих вычислялся как отношение числа побед бота к сумме побед и поражений бота. Результаты приведены в таблице~\ref{tab:baseline}.

\begin{table}[h!]
\captionsetup{position=above}
\caption{Результаты серии из 100 партий против игроков базового уровня}\label{tab:baseline}
\centering\small
\begin{tabular}{|>{\raggedright\arraybackslash}p{4.4cm}|c|c|c|c|>{\raggedright\arraybackslash}p{3.0cm}|}
\hline
Соперник & Побед & Пораж. & Ничьих & Ошибок & \% побед без ничьих\\ \hline
Базовый («лёгкий») & 94 & 4 & 2 & 0 & 95,9\\ \hline
Повышенный («сложный») & 25 & 2 & 73 & 0 & 92,6\\ \hline
\end{tabular}
\end{table}

Требование задания (более 50\% побед над игроком базового уровня и более 25\% над игроком повышенного уровня без учёта ничьих) выполнено с запасом: $94/(94+4)\approx 95{,}9\%$ и $25/(25+2)\approx 92{,}6\%$ соответственно. Ошибок и дисквалификаций ни в одной из 200 партий не было.

Против сильного соперника подавляющее большинство партий (73 из 100) заканчивается ничьей: оба игрока хорошо блокируют угрозы, и поле заполняется без выигрышной линии. Ничья не является поражением бота, поэтому результат полезно оценить и с учётом ничьих. Против лёгкого игрока бот не проиграл в $94+2=96$ партиях из 100 (96\%), против сильного --- в $25+73=98$ партиях из 100 (98\%). Доля побед от общего числа партий составила 94\% и 25\% соответственно. Критерий задания сформулирован без учёта ничьих, поэтому формальным показателем остаётся значение из последнего столбца таблицы~\ref{tab:baseline}.

\subsection{Время расчёта хода}

Время каждого хода измеряется программно: тест \code{test_stats_vs_baseline} выводит среднее время хода каждого игрока в миллисекундах (таблица~\ref{tab:time}).

\begin{table}[h!]
\captionsetup{position=above}
\caption{Время хода бота и длительность серий}\label{tab:time}
\centering\small
\begin{tabular}{|>{\raggedright\arraybackslash}p{6.4cm}|c|c|}
\hline
Характеристика & Против лёгкого & Против сильного\\ \hline
Среднее время хода бота, мс & 60,5 & 67,9\\ \hline
Среднее время хода соперника, мс & 0,002 & 0,003\\ \hline
Число партий в серии & 100 & 100\\ \hline
Общее время серии & 2 мин 43 с & 13 мин 8 с\\ \hline
\end{tabular}
\end{table}

Среднее время хода бота (60,5 и 67,9~мс) укладывается в требуемые 50--100~мс. Оно зависит от лимита 70~мс: на ходах, где перебор завершается раньше (например, при форсированных ходах), расчёт идёт быстрее, поэтому среднее значение ниже лимита. Максимальное время хода и достигнутая глубина перебора в данных тестах не фиксировались.

% ======================= ЗАКЛЮЧЕНИЕ =======================
\unnumsection{Заключение}

В рамках курсовой работы разработан бот-игрок в крестики-нолики до пяти в ряд на поле $20\times 20$ с препятствиями, реализующий поиск минимакс с альфа-бета отсечением, инкрементальной оценкой позиции по системе весов окон, отбором кандидатов на ход по радиусу и приоритету, форсированными ходами и итеративным углублением с ограничением по времени.

Реализация состоит из класса \code{MyPlayer} (около 320 строк кода в \code{my_player.cpp} и 72 строки в \code{my_player.hpp}) и набора из трёх дополнительных автоматических тестов на CTest (около 140 строк в \code{test_my_bot.cpp}), дополняющих уже имевшиеся в проекте тесты. Настроена конфигурация непрерывной интеграции GitHub Actions (\code{.github/workflows/ci.yml}), автоматически собирающая проект и запускающая тесты при каждом push и pull request в отслеживаемые ветки. Разработка велась через трекер задач: девять задач на Kanban-доске, каждая --- в отдельной ветке и с отдельным Pull Request в ветку \code{dev}.

Результаты испытаний: в серии из 100 партий бот выиграл у лёгкого игрока базового уровня 94 партии при 4 поражениях и 2 ничьих (95,9\% побед без учёта ничьих), у сильного --- 25 партий при 2 поражениях и 73 ничьих (92,6\%). Среднее время хода составило 60,5~мс против лёгкого и 67,9~мс против сильного игрока, что укладывается в требуемый диапазон 50--100~мс. \todo{Дополните: сколько времени суммарно заняли этапы работы --- согласование алгоритма, реализация, тестирование, оформление отчёта.}

Полученный алгоритм детерминирован, не требует генератора случайных чисел и работает по одному и тому же принципу вне зависимости от конкретной расстановки препятствий на поле, что упрощает его тестирование и объяснение.

Слабые места и возможные улучшения:
\begin{itemize}
  \item оценочная функция учитывает только окна длины ровно 5 и не различает, например, «открытую» тройку (с двух сторон от которой есть свободное место) от «закрытой» --- это можно уточнить, увеличив число градаций весов;
  \item правило «последнего хода» второго игрока после линии первого игрока не смоделировано отдельно в переборе --- учитывается только выигрыш, что является допустимым упрощением, но может занижать оценку позиций, близких к этому редкому случаю;
  \item ширина луча перебора (\code{beam}) подобрана эмпирически; более аккуратная настройка (например, адаптивная --- шире у корня, уже глубже) могла бы увеличить достигаемую глубину при том же лимите времени;
  \item не реализована таблица транспозиций (кэш уже просчитанных позиций), которая на играх с большим числом повторяющихся подпозиций обычно даёт существенный прирост эффективности перебора;
  \item тест \code{test_player_stats} (100 партий самоигры) занимает продолжительное время на отладочной сборке и в среде CI; в дальнейшем стоит рассмотреть выделение длительных статистических тестов в отдельный, запускаемый по требованию workflow, чтобы не замедлять проверку каждого коммита.
\end{itemize}

% ======================= СПИСОК ЛИТЕРАТУРЫ =======================
\begin{thebibliography}{9}
\renewcommand{\refname}{Список литературы}
\phantomsection\addcontentsline{toc}{section}{Список литературы}
\bibitem{russell} Рассел С., Норвиг П. Искусственный интеллект: современный подход. 3-е изд. М.: Вильямс, 2018.
\bibitem{repo} vsennov. tictactoe-course: материалы курсовой работы по программированию на C++ [Электронный ресурс] // GitHub. URL: \url{https://github.com/vsennov/tictactoe-course} (дата обращения: 25.10.2026).
\bibitem{cormen} Cormen T. H., Leiserson C. E., Rivest R. L., Stein C. Introduction to Algorithms. 3rd ed. Cambridge, MA: MIT Press, 2009.
\bibitem{cppref} C++ reference [Электронный ресурс] // cppreference.com. URL: \url{https://en.cppreference.com/} (дата обращения: 25.10.2026).
\bibitem{cmake} CMake Reference Documentation [Электронный ресурс]. URL: \url{https://cmake.org/cmake/help/latest/} (дата обращения: 25.10.2026).
\bibitem{ctest} ctest(1) --- CMake Documentation [Электронный ресурс]. URL: \url{https://cmake.org/cmake/help/latest/manual/ctest.1.html} (дата обращения: 25.10.2026).
\bibitem{ghactions} GitHub Actions documentation [Электронный ресурс] // GitHub Docs. URL: \url{https://docs.github.com/actions} (дата обращения: 25.10.2026).
\bibitem{ghprojects} Planning and tracking with Projects [Электронный ресурс] // GitHub Docs. URL: \url{https://docs.github.com/issues/planning-and-tracking-with-projects} (дата обращения: 25.10.2026).
\end{thebibliography}

% ======================= ПРИЛОЖЕНИЕ А =======================
\unnumsection{Приложение А. Исходный код бота, тестов и CI}

\subsection*{А.1. Файл \texttt{src/player/my\_player.hpp}}
\begin{verbatim}
#pragma once

#include "core/game.hpp"

#include <chrono>
#include <cstdint>
#include <vector>

namespace ttt::my_player {

using game::Event;
using game::IPlayer;
using game::Point;
using game::Sign;
using game::State;

// Бот: минимакс (негамакс) с альфа-бета отсечением, итеративным углублением,
// ограничением по времени и упорядочиванием ходов.
class MyPlayer : public IPlayer {
public:
  MyPlayer(const char *name) : m_sign(Sign::NONE), m_name(name) {}
  void set_sign(Sign sign) override;
  Point make_move(const State &game) override;
  const char *get_name() const override;

  // Настройка и статистика (для тестов и отчета).
  void set_time_limit_ms(int ms) { m_time_limit_ms = ms; }
  int last_depth() const { return m_last_depth; }
  long long last_nodes() const { return m_last_nodes; }
  double last_time_ms() const { return m_last_ms; }

private:
  struct Cand {
    int x;
    int y;
    long long pr;
  };

  Sign m_sign;
  const char *m_name;

  // Внутренняя доска: 0 - пусто, 1 - я, 2 - соперник, 3 - препятствие.
  int m_nx = 0, m_ny = 0, m_len = 5;
  std::vector<std::int8_t> m_b;
  std::vector<int> m_near; // сколько занятых клеток в радиусе 2
  long long m_score = 0;   // оценка позиции с точки зрения "меня"
  long long m_W[16], m_att[16], m_def[16];
  std::vector<std::vector<Cand>> m_cands;

  // Управление поиском.
  int m_time_limit_ms = 70;
  std::chrono::steady_clock::time_point m_deadline;
  bool m_timeout = false;
  long long m_nodes = 0;
  int m_last_depth = 0;
  long long m_last_nodes = 0;
  double m_last_ms = 0;

  void load(const State &game, Sign me);
  std::int8_t at(int x, int y) const;
  long long window_value(int sx, int sy, int d) const;
  void apply(int x, int y, std::int8_t p);
  void undo(int x, int y, long long saved);
  void bump_near(int x, int y, int delta);
  bool wins(int x, int y, std::int8_t p) const;
  long long cell_priority(int x, int y, std::int8_t mover) const;
  int gen(int ply, std::int8_t mover, int beam);
  int negamax(int depth, int ply, int alpha, int beta, std::int8_t side);
  int search_root(int depth, int hint_x, int hint_y, int &bx, int &by);
};

}; // namespace ttt::my_player
\end{verbatim}

\subsection*{А.2. Файл \texttt{src/player/my\_player.cpp}}
\begin{verbatim}
#include "player/my_player.hpp"

#include <algorithm>
#include <cstdlib>

namespace ttt::my_player {

namespace {
constexpr std::int8_t EMPTY = 0, ME = 1, OPP = 2, WALL = 3;
constexpr int DX[4] = {1, 0, 1, 1};
constexpr int DY[4] = {0, 1, 1, -1};
constexpr int INF = 1000000000;
constexpr int WIN = 100000000;
constexpr int ROOT_BEAM = 30; // сколько ходов рассматриваем в корне
constexpr int BEAM = 12;      // ... и на остальных уровнях
constexpr int MAX_DEPTH = 16;

long long pow10ll(int n) {
  long long r = 1;
  while (n-- > 0)
    r *= 10;
  return r;
}
std::int8_t other(std::int8_t s) { return s == ME ? OPP : ME; }
} // namespace

void MyPlayer::set_sign(Sign sign) { m_sign = sign; }

const char *MyPlayer::get_name() const { return m_name; }

std::int8_t MyPlayer::at(int x, int y) const {
  if (x < 0 || y < 0 || x >= m_nx || y >= m_ny)
    return WALL;
  return m_b[x * m_ny + y];
}

// Копируем поле из игры во внутреннее представление и считаем оценку.
void MyPlayer::load(const State &game, Sign me) {
  const auto &o = game.get_opts();
  m_nx = o.rows;
  m_ny = o.cols;
  m_len = std::min(std::max(o.win_len, 2), 15);
  m_b.assign(m_nx * m_ny, EMPTY);
  m_near.assign(m_nx * m_ny, 0);
  m_score = 0;
  m_cands.assign(MAX_DEPTH + 2, std::vector<Cand>());

  // Таблицы весов.
  m_W[0] = m_att[0] = m_def[0] = 0;
  for (int c = 1; c < m_len; ++c) {
    m_W[c] = pow10ll(c - 1);
    m_att[c] = pow10ll(c);
    m_def[c] = pow10ll(c) / 2;
  }
  if (m_len >= 4) {
    m_W[m_len - 1] = 20000; // четверка: следующим ходом выигрыш
    m_W[m_len - 2] = 300;   // тройка: угроза создать четверку
  }
  m_W[m_len] = 10000000;
  m_att[m_len - 1] = 10000000;
  m_def[m_len - 1] = 1000000;

  for (int x = 0; x < m_nx; ++x)
    for (int y = 0; y < m_ny; ++y) {
      Sign s = game.get_value(x, y);
      std::int8_t v = EMPTY;
      if (s == Sign::WALL)
        v = WALL;
      else if (s == Sign::NONE)
        v = EMPTY;
      else
        v = (s == me) ? ME : OPP;
      m_b[x * m_ny + y] = v;
    }
  for (int x = 0; x < m_nx; ++x)
    for (int y = 0; y < m_ny; ++y)
      if (m_b[x * m_ny + y] == ME || m_b[x * m_ny + y] == OPP)
        bump_near(x, y, +1);
  // Полная начальная оценка: сумма по всем окнам длины win_len.
  for (int x = 0; x < m_nx; ++x)
    for (int y = 0; y < m_ny; ++y)
      for (int d = 0; d < 4; ++d)
        m_score += window_value(x, y, d);
}

// Оценка одного окна из m_len клеток: >0 в пользу меня, <0 в пользу соперника.
long long MyPlayer::window_value(int sx, int sy, int d) const {
  int me = 0, opp = 0;
  for (int i = 0; i < m_len; ++i) {
    std::int8_t c = at(sx + DX[d] * i, sy + DY[d] * i);
    if (c == WALL)
      return 0;
    if (c == ME)
      ++me;
    else if (c == OPP)
      ++opp;
  }
  if (me && opp)
    return 0;
  if (me)
    return m_W[me];
  if (opp)
    return -m_W[opp];
  return 0;
}

void MyPlayer::bump_near(int x, int y, int delta) {
  for (int dx = -2; dx <= 2; ++dx)
    for (int dy = -2; dy <= 2; ++dy) {
      int nx = x + dx, ny = y + dy;
      if (nx >= 0 && ny >= 0 && nx < m_nx && ny < m_ny)
        m_near[nx * m_ny + ny] += delta;
    }
}

// Поставить фишку и пересчитать оценку только по окнам, содержащим клетку.
void MyPlayer::apply(int x, int y, std::int8_t p) {
  for (int d = 0; d < 4; ++d)
    for (int k = 0; k < m_len; ++k)
      m_score -= window_value(x - DX[d] * k, y - DY[d] * k, d);
  m_b[x * m_ny + y] = p;
  for (int d = 0; d < 4; ++d)
    for (int k = 0; k < m_len; ++k)
      m_score += window_value(x - DX[d] * k, y - DY[d] * k, d);
  bump_near(x, y, +1);
}

void MyPlayer::undo(int x, int y, long long saved) {
  m_b[x * m_ny + y] = EMPTY;
  bump_near(x, y, -1);
  m_score = saved;
}

// Собрана ли линия длины >= win_len через клетку (x, y)?
bool MyPlayer::wins(int x, int y, std::int8_t p) const {
  for (int d = 0; d < 4; ++d) {
    int cnt = 1;
    for (int s = 1; at(x + DX[d] * s, y + DY[d] * s) == p; ++s)
      ++cnt;
    for (int s = 1; at(x - DX[d] * s, y - DY[d] * s) == p; ++s)
      ++cnt;
    if (cnt >= m_len)
      return true;
  }
  return false;
}

// Быстрая оценка полезности хода в клетку (для упорядочивания ходов):
// атака (свои линии) + защита (блокировка линий соперника).
long long MyPlayer::cell_priority(int x, int y, std::int8_t mover) const {
  long long s = 0;
  for (int d = 0; d < 4; ++d)
    for (int k = 0; k < m_len; ++k) {
      int sx = x - DX[d] * k, sy = y - DY[d] * k;
      int mv = 0, ot = 0;
      bool bad = false;
      for (int i = 0; i < m_len; ++i) {
        std::int8_t c = at(sx + DX[d] * i, sy + DY[d] * i);
        if (c == WALL) {
          bad = true;
          break;
        }
        if (c == mover)
          ++mv;
        else if (c != EMPTY)
          ++ot;
      }
      if (bad || (mv && ot))
        continue;
      if (!mv && !ot)
        s += 1;
      else if (!ot)
        s += m_att[mv];
      else
        s += m_def[ot];
    }
  return s;
}

// Кандидаты: свободные клетки в радиусе 2 от занятых, лучшие beam по приоритету.
int MyPlayer::gen(int ply, std::int8_t mover, int beam) {
  auto &v = m_cands[ply];
  v.clear();
  for (int x = 0; x < m_nx; ++x)
    for (int y = 0; y < m_ny; ++y) {
      int idx = x * m_ny + y;
      if (m_b[idx] == EMPTY && m_near[idx] > 0)
        v.push_back({x, y, cell_priority(x, y, mover)});
    }
  int n = std::min<int>(beam, (int)v.size());
  std::partial_sort(v.begin(), v.begin() + n, v.end(),
                    [](const Cand &a, const Cand &b) { return a.pr > b.pr; });
  // Форсированные ходы: есть ход, собирающий линию, или надо блокировать
  // соперника, у которого до линии не хватает одной клетки.
  if (n > 0 && v[0].pr >= m_def[m_len - 1])
    n = 1;
  return n;
}

int MyPlayer::negamax(int depth, int ply, int alpha, int beta,
                      std::int8_t side) {
  if ((++m_nodes & 255) == 0 && std::chrono::steady_clock::now() >= m_deadline)
    m_timeout = true;
  if (m_timeout)
    return 0;
  int n = gen(ply, side, BEAM);
  if (n == 0)
    return 0; // нет ходов - считаем ничьей
  int best = -INF;
  for (int i = 0; i < n; ++i) {
    Cand c = m_cands[ply][i];
    long long saved = m_score;
    apply(c.x, c.y, side);
    int val;
    if (wins(c.x, c.y, side))
      val = WIN - ply;
    else if (depth <= 1)
      val = (int)(side == ME ? m_score : -m_score);
    else
      val = -negamax(depth - 1, ply + 1, -beta, -alpha, other(side));
    undo(c.x, c.y, saved);
    if (m_timeout)
      return 0;
    if (val > best)
      best = val;
    if (best > alpha)
      alpha = best;
    if (alpha >= beta)
      break;
  }
  return best;
}

// Один проход итеративного углубления; возвращает оценку лучшего хода.
int MyPlayer::search_root(int depth, int hint_x, int hint_y, int &bx,
                          int &by) {
  int n = gen(0, ME, ROOT_BEAM);
  auto &v = m_cands[0];
  for (int i = 0; i < n; ++i) // ход с прошлой итерации ставим первым
    if (v[i].x == hint_x && v[i].y == hint_y) {
      std::rotate(v.begin(), v.begin() + i, v.begin() + i + 1);
      break;
    }
  int alpha = -INF, best = -INF;
  bx = by = -1;
  for (int i = 0; i < n; ++i) {
    Cand c = v[i];
    long long saved = m_score;
    apply(c.x, c.y, ME);
    int val;
    if (wins(c.x, c.y, ME))
      val = WIN;
    else
      val = -negamax(depth - 1, 1, -INF, -alpha, OPP);
    undo(c.x, c.y, saved);
    if (m_timeout)
      return 0;
    if (val > best) {
      best = val;
      bx = c.x;
      by = c.y;
    }
    if (best > alpha)
      alpha = best;
  }
  return best;
}

Point MyPlayer::make_move(const State &game) {
  auto t0 = std::chrono::steady_clock::now();
  load(game, game.get_current_player());
  m_deadline = t0 + std::chrono::milliseconds(m_time_limit_ms);
  m_timeout = false;
  m_nodes = 0;
  m_last_depth = 0;

  int bx = -1, by = -1;
  int n = gen(0, ME, ROOT_BEAM);
  if (n > 0) {
    bx = m_cands[0][0].x; // запасной ход - лучший по приоритету
    by = m_cands[0][0].y;
  } else { // на поле нет фишек: ближайшая к центру свободная клетка
    int bd = INF;
    for (int x = 0; x < m_nx; ++x)
      for (int y = 0; y < m_ny; ++y)
        if (m_b[x * m_ny + y] == EMPTY) {
          int d = std::abs(2 * x - (m_nx - 1)) + std::abs(2 * y - (m_ny - 1));
          if (d < bd) {
            bd = d;
            bx = x;
            by = y;
          }
        }
  }

  if (n > 0) {
    for (int depth = 2; depth <= MAX_DEPTH; depth += 2) {
      int rx, ry;
      int val = search_root(depth, bx, by, rx, ry);
      if (m_timeout || rx < 0)
        break;
      bx = rx;
      by = ry;
      m_last_depth = depth;
      if (std::abs(val) >= WIN - 100)
        break; // найден форсированный выигрыш/проигрыш
    }
  }

  m_last_nodes = m_nodes;
  m_last_ms = std::chrono::duration<double, std::milli>(
                  std::chrono::steady_clock::now() - t0)
                  .count();
  if (bx < 0)
    return Point{0, 0};
  return Point{bx, by};
}

}; // namespace ttt::my_player
\end{verbatim}

\subsection*{А.3. Файл \texttt{tests/test\_my\_bot.cpp}}
\begin{verbatim}
// Автоматические тесты для MyPlayer (CTest).
// Файл кладётся в tests/, рядом с остальными test_*.cpp.
//
// Если в вашем test/CMakeLists.txt тесты добавляются вручную, добавьте:
//   add_executable(test_my_bot test_my_bot.cpp)
//   target_link_libraries(test_my_bot PRIVATE tttcore tttplayer)
//   add_test(NAME test_my_bot COMMAND test_my_bot)
// Если используется макрос/функция для регистрации тестов (см. другие
// test_*.cpp в проекте), зарегистрируйте test_my_bot так же, как соседние.

#include "core/game.hpp"
#include "player/my_player.hpp"

#include <cassert>
#include <chrono>
#include <cstdio>
#include <cstdlib>

using namespace ttt::game;
using ttt::my_player::MyPlayer;

namespace {

State::Opts make_opts() {
  State::Opts o;
  o.rows = o.cols = 20;
  o.win_len = 5;
  o.max_moves = 0;
  return o;
}

// Тест 1: бот всегда возвращает координаты в пределах поля и в свободную
// клетку — на пустом поле и на поле с уже сделанными ходами.
bool test_move_is_legal() {
  auto opts = make_opts();
  RandomObstaclesFI fi(0.75, 50, 1);
  State state(opts, &fi);
  MyPlayer bot("bot");
  bot.set_sign(Sign::X);

  for (int i = 0; i < 5; ++i) {
    Point p = bot.make_move(state);
    if (p.x < 0 || p.x >= opts.rows || p.y < 0 || p.y >= opts.cols) {
      std::printf("test_move_is_legal: координаты вне поля (%d, %d)\n", p.x,
                  p.y);
      return false;
    }
    if (state.get_value(p.x, p.y) != Sign::NONE) {
      std::printf("test_move_is_legal: клетка (%d, %d) занята\n", p.x, p.y);
      return false;
    }
    MoveResult r = state.process_move(Sign::X, p.x, p.y);
    if (r != MoveResult::OK && r != MoveResult::WIN && r != MoveResult::DRAW) {
      std::printf("test_move_is_legal: неожиданный результат хода\n");
      return false;
    }
    if (r != MoveResult::OK)
      break;
    // Ход соперника случайной свободной клеткой, чтобы игра продолжалась.
    bool moved = false;
    for (int x = 0; x < opts.rows && !moved; ++x)
      for (int y = 0; y < opts.cols && !moved; ++y)
        if (state.get_value(x, y) == Sign::NONE) {
          state.process_move(Sign::O, x, y);
          moved = true;
        }
  }
  return true;
}

// Тест 2: полная партия бота против самого себя (два экземпляра) не
// заканчивается дисквалификацией ни одной из сторон.
bool test_self_play_no_dq() {
  auto opts = make_opts();
  RandomObstaclesFI fi(0.75, 50, 1);
  State state(opts, &fi);

  MyPlayer px("botX");
  MyPlayer po("botO");
  px.set_sign(Sign::X);
  po.set_sign(Sign::O);
  px.set_time_limit_ms(30);
  po.set_time_limit_ms(30);

  Game game(opts, &fi);
  game.add_player(Sign::X, &px);
  game.add_player(Sign::O, &po);

  MoveResult r;
  int guard = 0;
  do {
    r = game.process();
    ++guard;
  } while (r == MoveResult::OK && guard < opts.rows * opts.cols + 5);

  if (is_dq(r)) {
    std::printf("test_self_play_no_dq: дисквалификация, result=%d\n",
               (int)r);
    return false;
  }
  return true;
}

// Тест 3: время одного хода укладывается в установленный лимит с запасом.
bool test_move_time_limit() {
  auto opts = make_opts();
  RandomObstaclesFI fi(0.75, 50, 1);
  State state(opts, &fi);
  MyPlayer bot("bot");
  bot.set_sign(Sign::X);
  bot.set_time_limit_ms(70);

  auto t0 = std::chrono::steady_clock::now();
  bot.make_move(state);
  double ms = std::chrono::duration<double, std::milli>(
                  std::chrono::steady_clock::now() - t0)
                  .count();

  if (ms > 100.0) {
    std::printf("test_move_time_limit: ход занял %.1f мс (> 100 мс)\n", ms);
    return false;
  }
  return true;
}

} // namespace

int main() {
  struct { const char *name; bool (*fn)(); } tests[] = {
      {"test_move_is_legal", test_move_is_legal},
      {"test_self_play_no_dq", test_self_play_no_dq},
      {"test_move_time_limit", test_move_time_limit},
  };
  bool all_ok = true;
  for (auto &t : tests) {
    bool ok = t.fn();
    std::printf("%s: %s\n", t.name, ok ? "OK" : "FAIL");
    all_ok = all_ok && ok;
  }
  return all_ok ? 0 : 1;
}
\end{verbatim}

\subsection*{А.4. Файл \texttt{.github/workflows/ci.yml}}
\begin{verbatim}
# Кладётся в .github/workflows/ci.yml
# Собирает проект (без предкомпилированной библиотеки baseline, так как она
# закрытая и недоступна CI) и запускает CTest.
name: CI

on:
  push:
    branches: [ main, dev, "feature/**" ]
  pull_request:
    branches: [ main, dev ]

jobs:
  build-and-test:
    runs-on: ubuntu-latest
    steps:
      - name: Checkout
        uses: actions/checkout@v4

      - name: Configure
        run: |
          mkdir -p build
          cd build
          cmake .. -DCMAKE_BUILD_TYPE=Release

      - name: Build
        run: cmake --build build --config Release -j

      - name: Test
        run: |
          cd build
          ctest --output-on-failure
\end{verbatim}

\end{document}
